\begin{trapbox}
	\begin{description}[style=nextline, leftmargin=2.8em, labelsep=0.5em, itemsep=0.35em]
		\item[\textbf{\textcolor{CTrap}{易错点一（符号混淆）}}] 误将左焦半径记作 $a-ex$，右焦半径记作 $a+ex$。
		\item[\textbf{\textcolor{CTrap}{正确理解（左加右减）：}}] 对于焦点
		      \[
			      F_1(-c,0),\qquad F_2(c,0),
		      \]
		      有
		      \[
			      PF_1=a+ex,\qquad PF_2=a-ex.
		      \]
		      口诀：“左加右减”。
		\item[\textbf{\textcolor{CTrap}{错因分析（焦点位置混淆）：}}] 忘记 $x$ 前的符号与焦点坐标符号有关，机械记忆导致颠倒。
	\end{description}

	\medskip
	\begin{description}[style=nextline, leftmargin=2.8em, labelsep=0.5em, itemsep=0.35em]
		\item[\textbf{\textcolor{CTrap}{易错点二（忽视 $x$ 范围）}}] 利用公式 $PF_2=a-ex$ 求 $x$ 时，未检查 $x$ 是否在 $[-a,a]$ 内，导致得出点不在椭圆上的错误答案。
		\item[\textbf{\textcolor{CTrap}{正确理解（检验定义域）：}}] 解出 $x$ 后必须验证
		      \[
			      |x|\le a.
		      \]
		      否则该点不在椭圆上，应舍去。
		\item[\textbf{\textcolor{CTrap}{进一步检查（焦半径范围）：}}] 因为
		      \[
			      -a\le x\le a,\qquad e=\frac{c}{a},
		      \]
		      所以
		      \[
			      a-c\le PF_1\le a+c,\qquad a-c\le PF_2\le a+c.
		      \]
		      若题目给出的焦半径不在 $[a-c,a+c]$ 内，则这样的点 $P$ 不存在。
		\item[\textbf{\textcolor{CTrap}{错因分析（代数解强制接受）：}}] 只解方程不检验几何约束，将椭圆方程的取值范围忽略。
	\end{description}

	\medskip
	\begin{description}[style=nextline, leftmargin=2.8em, labelsep=0.5em, itemsep=0.35em]
		\item[\textbf{\textcolor{CTrap}{易错点三（把有向差当成绝对值）}}] 看到
		      \[
			      PF_1-PF_2=2ex
		      \]
		      后，误以为“两焦半径之差”一定等于 $2ex$ 且一定非负。
		\item[\textbf{\textcolor{CTrap}{正确理解（有向差与绝对差）：}}]
		      \[
			      PF_1-PF_2=2ex
		      \]
		      是有向差，可以为正、为零或为负。若题目问“两焦半径之差的绝对值”，则应写成
		      \[
			      |PF_1-PF_2|=2e|x|.
		      \]
		\item[\textbf{\textcolor{CTrap}{错因分析（忽视 $x$ 的正负）：}}] 当 $x<0$ 时，$PF_1-PF_2<0$，这并不矛盾，因为它表示的是有向差。
	\end{description}

	\medskip
	\begin{description}[style=nextline, leftmargin=2.8em, labelsep=0.5em, itemsep=0.35em]
		\item[\textbf{\textcolor{CTrap}{易错点四（公式前提遗忘）}}] 不考虑椭圆焦点位置，直接套用 $PF=a\pm ex$ 到焦点在 $y$ 轴的椭圆。
		\item[\textbf{\textcolor{CTrap}{正确理解（焦点轴判断）：}}] 焦半径公式
		      \[
			      PF_1=a+ex,\qquad PF_2=a-ex
		      \]
		      针对焦点在 $x$ 轴，且
		      \[
			      F_1(-c,0),\qquad F_2(c,0)
		      \]
		      的标准椭圆。

		      若椭圆为
		      \[
			      \frac{x^2}{b^2}+\frac{y^2}{a^2}=1\quad (a>b>0),
		      \]
		      焦点为
		      \[
			      F_1(0,-c),\qquad F_2(0,c),
		      \]
		      则
		      \[
			      PF_1=a+ey,\qquad PF_2=a-ey.
		      \]
		      也就是说：负方向焦点用加号，正方向焦点用减号。若焦点命名反过来，公式中的符号也要随之改变。
		\item[\textbf{\textcolor{CTrap}{错因分析（标准方程误用）：}}] 未注意椭圆标准方程中 $a$ 对应的轴，也未注意焦点的命名顺序，直接机械套用结论。
	\end{description}
\end{trapbox}